\chapter{Results: Terminal Structural Condition}\label{ch:structural}
This chapter answers RQ2--RQ4 for the measured structural endpoints. It first distinguishes robust model selection from diagnostic winners, then concentrates on the strongest controlled question: the incremental value of image features for terminal load.

\section{Two measured endpoints, weaker transfer}
In the earlier interpretable phase, grouped-holdout RandomForest gives wire-area-loss MAE 3.899 percentage points and $R^2=0.310$, while ultimate-load MAE is 0.121\,kN and $R^2=0.385$. These scores use ten test specimens from 48, with 38 training specimens. They show some within-mixture predictability, but do not establish stable transfer.

The later robustness-weighted selection chooses metadata-only models for both endpoints (Table~\ref{tab:robustselection}). For wire-area loss the final selection is CatBoost, with grouped MAE 0.124 as a fraction. A separate diagnostic table follows a RandomForest with grouped MAE 0.105. Those rows must not be merged: the lower diagnostic number is not the final CatBoost model's score.
\begin{table}[htbp]\centering\small
\caption{Final robustness-selected structural configurations on the 791-image basis with 48 terminal specimens. Both use metadata only. Entries are mean MAE across five grouped holdouts, nine treatment folds, or two campaign folds; wire-loss error is a fraction, and load error is in kN. Selection weighs evidence across regimes, so these are not necessarily the smallest isolated grouped errors.}\label{tab:robustselection}
\begin{tabular}{llrrr}\toprule
Target & Model & Grouped & LOTO & LOCO\\\midrule
Wire-area loss (fraction) & CatBoost & 0.124 & 0.116 & 0.121\\
Ultimate load (kN) & RandomForest & 0.174 & 0.203 & 0.564\\\bottomrule
\end{tabular}
\end{table}
The final load model's campaign-out MAE is 3.236 times its own grouped MAE, and its mean rank correlation falls from 0.778 to 0.353. The diagnostic wire-loss RandomForest has campaign-out Spearman $-0.089$. This combination of modest within-mixture error and poor transfer motivates the focused load study, rather than a direct image-to-life claim.

\section{Do images improve the metadata baseline?}\label{sec:ablation}
Table~\ref{tab:ablation} gives all seven selected feature families in the pooled terminal-load experiment. Metrics use terminal held-out observations, with 48 specimens covered in each of two five-fold repetitions. The metadata-only Ridge configuration has mean MAE 0.173\,kN, fold SD 0.046\,kN, mean RMSE 0.214\,kN, and mean Spearman 0.758.
\begin{table}[htbp]\centering\small
\caption{Pooled terminal-load comparison for 48 specimens across ten grouped folds. MAE and its sample SD are in kN; $R^2$ and Spearman $\rho$ are arithmetic means of fold metrics. Each feature family shows its selected model, chosen using the same saved folds. The table is an exploratory comparison, not a nested estimate of the model-selection procedure or a significance test.}\label{tab:ablation}
\input{\shared/tables/pooled_capacity.tex}
\end{table}
Adding HSV descriptors to metadata changes mean MAE only from 0.172711 to 0.173003\,kN while increasing mean rank correlation from 0.758 to 0.771. The MAE change is about 0.000293\,kN. It is not evidence of an improvement, and the available comparison does not establish equivalence either. Image-only feature families have larger pooled MAE than the metadata anchor. The combined RGB+HSV configuration with metadata is also higher in MAE than the anchor.

\fig{capacity_ablation}{1}{Pooled terminal-load results for selected feature-family models on 48 independent specimens across ten grouped folds. Points are arithmetic means and horizontal bars are sample standard deviations across folds, not confidence intervals. Terminal test sets contain nine or ten specimens. All training photographs share their specimen's terminal target; scoring is at terminal observations. Models are identified in Table~\ref{tab:ablation}. The figure summarizes the same saved comparison, not an additional experiment.}{fig:ablation}

The observed pattern supports a bounded conclusion: the tested image descriptors do not provide a reliable pooled increment over metadata for terminal load in this experiment. It does not show that images are irrelevant to structural behaviour. It also does not show that metadata effects are causal or transferable outside the experimental mixture.

\section{The limited post-onset result}
The post-onset sensitivity analysis retains 43 terminal specimens after the first total-rust observation of at least 0.5\%. Its metadata+RGB+HSV CatBoost configuration has mean MAE 0.163\,kN and $R^2=0.660$, compared with metadata-only Ridge MAE 0.173\,kN and $R^2=0.633$. Mean Spearman is 0.791 for the combined model and 0.803 for metadata alone. Thus the observed improvement is limited to some metrics, rather than a uniform increase in predictive quality.

The selected model family changes as well as the input features, and the subset is smaller than the primary pooled basis. This is an exploratory sensitivity result. It motivates a future matched-model, predeclared study of image availability after onset; it does not overturn the pooled finding or establish a clinically or mechanically meaningful gain.

\section{Design stratification and campaign exclusion}\label{sec:mesh}
Figure~\ref{fig:confounding} combines the campaign design with specimen-level terminal associations. Total surface rust and load have pooled Spearman correlation 0.462 for 48 specimens. Within the 24-specimen mesh groups, the corresponding correlations are $-0.046$ for four meshes and $-0.401$ for seven meshes. The first is approximately null; the second is negative. This demonstrates the sensitivity of the pooled association to the group structure without decomposing its causal components.

\fig{confounding}{1}{Campaign design and terminal surface-rust/load associations. The table shows the two observed combinations of mesh, chloride, and terminal exposure. Scatter plots use one terminal observation per specimen: pooled $n=48$, four-mesh $n=24$, seven-mesh $n=24$. Orange points are four-mesh specimens; blue points are seven-mesh specimens. Correlations are independently recomputed from the saved specimen summary. The pooled positive association does not establish an independent effect of rust on load.}{fig:confounding}

All selected feature-family configurations have negative mean test $R^2$ within either mesh family. In the four-mesh group, values range from $-3.557$ to $-0.806$; in the seven-mesh group they range from $-4.654$ to $-0.298$. Absolute MAE can still look moderate because the target range is narrower. For example, the selected seven-mesh HSV model has MAE 0.137\,kN and Spearman 0.460, but mean $R^2=-0.386$. Such a result does not justify the claim that every subgroup model contains no ranking information; it does limit a claim of reliably accurate capacity prediction. Full subgroup results are retained in Appendix~\ref{app:reproduction}.

Under campaign exclusion, the best metadata-only family uses GradientBoosting, with mean MAE 0.465\,kN and mean $R^2=-4.639$ over the two directions. This differs from the pooled Ridge selection, so dividing these two scores would not be a fixed-model robustness ratio. The result shows poor performance even after choosing the best metadata-only model within this stress-test evaluation.

\section{Residual correction and structural conclusion}
A specimen-summary refinement tests whether corrosion summaries can correct metadata-prediction residuals within mesh families. Its selected pipeline applies no residual correction: stage-one Ridge alone gives mean MAE approximately 0.173\,kN and mean Spearman 0.749. The attempted correction models do not improve that baseline in the saved comparison. This is a separate corroborating analysis within the same dataset, not external validation.

The structural evidence converges on a limited capability: pooled prediction can exploit specimen and exposure metadata, while incremental image value and transfer are not established reliably. Sparse terminal supervision and campaign confounding are consistent with the weakness, but the experiments do not isolate their relative contribution from representation limitations or unmeasured variation. The following chapter examines what happens when these structural estimates are propagated into time-dependent screening.
